\documentclass[10pt,A4]{article}
\usepackage[utf8]{inputenc}
\usepackage{xifthen}
\usepackage[default]{gillius}
\renewcommand*\familydefault{\sfdefault}
\usepackage[T1]{fontenc}
\usepackage{moresize}
\usepackage{fontawesome}
\usepackage[a4paper]{geometry}
\geometry{top=1cm, bottom=-.6cm, left=0.4cm, right=1cm}
\setlength{\headheight}{-5pt}
\setlength{\parindent}{0mm}
\usepackage{multicol}
\usepackage{multirow}
\usepackage{array}
\newcolumntype{x}[1]{%
>{\raggedleft\hspace{0pt}}p{#1}}%
\usepackage{graphicx}
\usepackage{wrapfig}
\usepackage{float}
\usepackage{tikz}
\usetikzlibrary{shapes, backgrounds,mindmap, trees}
\usepackage{transparent}
\usepackage{color}

% main colors
\definecolor{complcol}{RGB}{244,160,0}
\definecolor{bgcol}{RGB}{110,110,110}

\definecolor{softcol}{RGB}{219,68,55}
\definecolor{sectcol}{RGB}{66,133,244}
\definecolor{maya}{RGB}{124,185,232}
\definecolor{darkcol}{RGB}{51,51,51}
\definecolor{bgreycol}{RGB}{153,153,153}
\definecolor{fgreycol}{RGB}{85,85,85}
\definecolor{lightcol}{RGB}{238,238,221}

% pie chart colors bgcol:gris
% softcol:rojo
% sectcol:azul fuerte
% complcol: amarillo
% maya: azul bajo

\definecolor{col1}{RGB}{66,133,244}
\definecolor{col2}{RGB}{219,68,55}
\definecolor{col3}{RGB}{244,160,0}
\definecolor{col4}{RGB}{15,157,88}
\definecolor{col5}{RGB}{251,120,56}

\usepackage[hidelinks]{hyperref}

\newcommand{\mpwidth}{\linewidth-\fboxsep-\fboxsep}

\newcommand{\tzlarrow}{(0,0) -- (0.2,0) -- (0.3,0.2) -- (0.2,0.4) -- (0,0.4) -- (0.1,0.2) -- cycle;}

% include the left arrow into a tikz picture
% param1: fill color
%
\newcommand{\larrow}[1]
{\begin{tikzpicture}[scale=0.58]
	 \filldraw[fill=#1!100,draw=#1!100!black]  \tzlarrow
 \end{tikzpicture}
}

% a six pointed arrow poiting to the right
\newcommand{\tzrarrow}{ (0,0.2) -- (0.1,0) -- (0.3,0) -- (0.2,0.2) -- (0.3,0.4) -- (0.1,0.4) -- cycle;}

% include the right arrow into a tikz picture
% param1: fill color
%
\newcommand{\rarrow}
{\begin{tikzpicture}[scale=0.7]
	\filldraw[fill=softcol!100,draw=softcol!100!black] \tzrarrow
 \end{tikzpicture}
}

%----------------------------------------------------------------------------------------
%	custom sections
%----------------------------------------------------------------------------------------

% create a coloured box with arrow and title as cv section headline
% param 1: section title
%
\newcommand{\cvsection}[1]
{
\colorbox{fgreycol}{\mystrut \makebox[1\mpwidth][l]{
\larrow{softcol} \hspace{-8pt} \larrow{softcol} \hspace{-8pt} \larrow{softcol} \textbf{\textcolor{lightcol}{\uppercase{#1}}}\hspace{4pt}
}}\\
}

\newcommand{\cvskill}[3] {
	\begin{tabular*}{1\mpwidth}{p{0.75\mpwidth}  r}
		\textcolor{lightcol}{\textbf{#1}} & \textcolor{lightcol}{\textbf{#2}}\\
	\end{tabular*}%
	\hspace{1pt}
	\begin{tikzpicture}[scale=1,rounded corners=2pt,very thin]
	\fill [lightcol] (0,0) rectangle (1\mpwidth, 0.15);
	\fill [col4] (0,0) rectangle (#3\mpwidth, 0.15);
	\end{tikzpicture}%
}

% create a coloured arrow with title as cv meta section section
% param 1: meta section title
%
\newenvironment{metasection}[1] {
	\vspace{6pt}
	\begin{center}
		\textcolor{maya}{\large{\uppercase{#1}}}\\
	\normalsize
	\parbox{0.7\mpwidth}{\textcolor{maya}	\hrule}
}{\end{center}}

%----------------------------------------------------------------------------------------
%	 CV EVENT
%----------------------------------------------------------------------------------------

% creates a stretched box as cv entry headline followed by two paragraphs about
% the work you did
% param 1:	event time i.e. 2014 or 2011-2014 etc.
% param 2:	event name (what did you do?)
% param 3:	institution (where did you work / study)
% param 4:	what was your position
% param 5:	some words about your contributions
%
\newcommand{\cvevent}[5]
{
\vspace{8pt}
	\begin{tabular*}{1\mpwidth}{p{0.55\mpwidth}  x{0.42\mpwidth}}
 	\textcolor{black}{\textbf{#2}} & \textcolor{softcol}{#3}, \textcolor{darkcol}{\textbf{#1}}

	\end{tabular*}
\vspace{-12pt}
\textcolor{softcol}{\hrule}
\vspace{6pt}
	\begin{tabular*}{0.5\mpwidth}{p{\mpwidth}}
\larrow{softcol}  #4\\[3pt]
\larrow{softcol}  #5\\[6pt]
	\end{tabular*}

}

% creates a stretched box as
\newcommand{\cveventmeta}[2]
{
	\mbox{\mystrut \hspace{87pt}\textit{#1}}\\
	#2
}

%----------------------------------------------------------------------------------------
% CUSTOM STRUT FOR EMPTY BOXES
%----------------------------------------- -----------------------------------------------
\newcommand{\mystrut}{\rule[-.3\baselineskip]{0pt}{\baselineskip}}

%----------------------------------------------------------------------------------------
% CUSTOM LOREM IPSUM
%----------------------------------------------------------------------------------------
\newcommand{\lorem}
{Lorem ipsum dolor sit amet, consectetur adipiscing elit. Donec a diam lectus.}


% use to vertically center content
% credits to: http://tex.stackexchange.com/questions/7219/how-to-vertically-center-two-images-next-to-each-other
\newcommand{\vcenteredinclude}[1]{\begingroup
\setbox0=\hbox{\includegraphics{#1}}%
\parbox{\wd0}{\box0}\endgroup}

% use to vertically center content
% credits to: http://tex.stackexchange.com/questions/7219/how-to-vertically-center-two-images-next-to-each-other
\newcommand*{\vcenteredhbox}[1]{\begingroup
\setbox0=\hbox{#1}\parbox{\wd0}{\box0}\endgroup}

%----------------------------------------------------------------------------------------
%	ICON-SET EMBEDDING
%----------------------------------------------------------------------------------------

% at this point we simplify our icon-embedding by simply referring to a set of png images.
% if you find a good way of including svg without conflicting with other packages you can
% replace this part
\newcommand{\icon}[3]{\makebox(#2, #2){\textcolor{#3}{\csname fa#1\endcsname}}}	%icon shortcut
\newcommand{\icontext}[4]{ 						%icon with text shortcut
	\vcenteredhbox{\icon{#1}{#2}{#4}} \vcenteredhbox{\textcolor{#4}{#3}}
}
\newcommand{\iconhref}[5]{ 						%icon with website url
    \vcenteredhbox{\icon{#1}{#2}{#5}} \href{#4}{\textcolor{#5}{#3}}
}

\newcommand{\iconemail}[5]{ 						%icon with email link
    \vcenteredhbox{\icon{#1}{#2}{#5}} \href{mailto:#4}{\textcolor{#5}{#3}}
}

%----------------------------------------------------------------------------------------
% 	PIE CHART
%----------------------------------------------------------------------------------------
\input{./docs/pie-chart.tex}
